% =====================================================================
%  Tóm tắt
% =====================================================================
\section*{Tóm tắt}
\addcontentsline{toc}{section}{Tóm tắt}

\RIncomplete

\noindent
Báo cáo này trình bày một quy trình thực nghiệm có thể tái lập để đánh giá ảnh hưởng của
một \emph{nhiệm vụ phụ tự giám sát} đối với bài toán phân loại ảnh khi ngân sách nhãn bị
giới hạn. Cụ thể, CIFAR-10 được gộp thành hai siêu lớp do đề tài quy ước là
\emph{Animal} (6 lớp gốc) và \emph{Vehicle} (4 lớp gốc); một encoder ResNet-18 điều chỉnh cho
ảnh $32 \times 32$ được dùng chung cho hai đầu ra: phân loại siêu lớp có giám sát và dự đoán
góc xoay $k \in \{0,1,2,3\}$ được sinh tự động bằng \code{torch.rot90}. Nhãn của nhiệm vụ
phụ không cần nhãn ngữ nghĩa, nên nhánh xoay có thể khai thác thêm phần ảnh đã bị ẩn nhãn
của tập huấn luyện.

Khối thực nghiệm chính gồm ba cấu hình (\cfg{E1} baseline chỉ học siêu lớp, \cfg{E2-L} thêm
nhánh xoay nhưng chỉ dùng ảnh có nhãn, \cfg{E2} thêm nhánh xoay và dùng toàn bộ $\Dtrain$),
ba mức tỷ lệ nhãn ($10\,\%$, $20\,\%$, $50\,\%$) và ba lần lặp ghép cặp, tổng cộng \RRuns{}
lượt huấn luyện, mỗi lượt 100 epoch với tổng thời gian huấn luyện \RTotalHours{} giờ.
Chỉ số chính là Macro-F1 trên toàn bộ 10.000 ảnh test; Accuracy, precision/recall/F1 từng lớp
và ma trận nhầm lẫn được báo cáo kèm theo.

\medskip
\noindent\textbf{Kết quả chính.}
Tại mức nhãn $10\,\%$, Macro-F1 trên tập test đạt \REoneTenVal{} cho \cfg{E1} và \REtwoTenVal{}
cho \cfg{E2} (chênh lệch ghép cặp trung bình \REtwoTenDelta{} điểm phần trăm); tại mức
$50\,\%$ nhãn, \cfg{E1} đạt \REoneFiftyVal{} và \cfg{E2} đạt \REtwoFiftyVal{}
(chênh lệch \REtwoFiftyDelta{} điểm phần trăm). \RFindingSentence

\medskip
\noindent\textbf{Đóng góp có thể kiểm chứng.}
(i)~Một phép chia benchmark được đóng băng (seed 2026) với các tập nhãn lồng nhau
$\DL^{(10\%)} \subset \DL^{(20\%)} \subset \DL^{(50\%)}$ và chỉ mục đầy đủ được lưu lại;
(ii)~một giao thức kiểm tra việc che nhãn ở mức chỉ mục, bảo đảm nhãn thật của $\DU$ không
đi vào loss, không dùng để chọn positive pair và không tham gia chọn checkpoint;
(iii)~một khối kết quả ghép cặp theo lần lặp với chênh lệch $\Delta_s$ và độ lệch chuẩn
tường minh, kèm phân tích trường hợp không cải thiện;
(iv)~cài đặt đầy đủ \emph{bốn} cấu hình mở rộng contrastive \cfg{E3}--\cfg{E6} (NT-Xent và
SupCon) với siêu tham số $\lambda_c$ được khóa \emph{chỉ bằng validation} qua một quy trình
hai bước tách bạch thông tin;
(v)~các phép đo chi phí huấn luyện và suy luận được tách riêng, với kết luận về IoT được
giới hạn đúng theo thiết bị đã đo.

\medskip
\noindent\textbf{Kết quả phần mở rộng contrastive.}
\RExtensionConclusion

\medskip
\noindent\textbf{Giới hạn.}
Kết luận chỉ áp dụng cho phép gộp hai siêu lớp, các tỷ lệ nhãn và giao thức đã khảo sát.
Ba lần lặp chỉ cung cấp bằng chứng ban đầu về độ ổn định, không đủ để khẳng định ý nghĩa
thống kê. Kết quả đo độ trễ trên CPU của máy tính dùng để huấn luyện không phải là bằng chứng
về khả năng triển khai trên vi điều khiển. Không có kết quả nào trong báo cáo này được suy ra
từ số liệu minh họa.
